\section{Conclusion}

Selecting smoothness for a trend that
will be extrapolated is a forecasting
problem, not merely a post-hoc
descriptive smoothing problem.
We formulated chronological,
horizon-matched forecast
cross-validation over a normalized
PLS smoothness index $S\in[0,1]$.
The criterion compares historical
future-block errors from an explicitly
defined continuation, includes the
two limiting smoothing regimes, and
refits the selected trend before
issuing a genuinely out-of-sample
forecast.

The PLS smoother, trace-based
smoothness index, and predictive
cross-validation principle are
established methods. The contribution
studied here is their particular
operational combination for
finite-difference trend extrapolation
and a frozen simulation comparison
with conventional smoothing criteria.
At horizons $3$, $6$, and $12$,
the simulated geometric RMSFE
ratios to one-step tuning were
$0.947$, $0.861$, and $0.762$,
respectively, supporting the value
of selecting smoothness for the
forecast horizon within this design.

The objective can have multiple
local minima, and changing from
$\lambda$ to $S$ does not remove them.
Consequently, practical selection
must compare competing numerical
candidates and both endpoints.
The evidence does not establish
universal superiority of forecast-CV,
optimality across all polynomial
continuation orders, or a general
gain from tracked-minimum adaptation.
The principal result is a transparent
prediction-targeted tuning criterion
whose modeling assumptions,
computational search, and
information set can be stated
and evaluated explicitly.
